\section{Bayesian model surgery needs a separate prediction contract}
\label{exec23:pymc-scope}

PyMC is a plausible worker for bank-specific probabilistic models, including
hierarchical likelihoods and explicitly modeled latent quantities. The proposed
product, however, needs to distinguish three questions: which model is being
conditioned on data, which mechanism is replaced by an intervention, and which
variables are subsequently regenerated. The inspected repository revision is
dated 7 October 2026 and is identified by its immutable commit and twelve
authenticated source, documentation and test files. Its version is dynamically
resolved; this review does not assert that the selected development revision is
a named stable release. Its metadata requires Python at least 3.12 and below
3.15, and a specified PyTensor compatibility range. Those requirements, the
root Apache license and the actual dependency artifacts must be part of a
future worker contract rather than inferred from an older installed environment.
\citep{exec23_pymc_sources}

\subsection{Conditioning does not replace a treatment mechanism}
\label{exec23:pymc-intervention}

The inspected \texttt{observe} transformation converts selected random or
deterministic variables into observed variables, with compatible shape and
type. The \texttt{do} transformation instead replaces named model variables
by intervention expressions in a cloned graph. An expression depending on
random variables is represented as a deterministic; a constant intervention
is shared data by default. These operations provide useful model construction
machinery. They do not establish that the supplied graph correctly represents
assignment, latent common causes, information timing or a bank's feasible
interventions.
\citep{exec23_pymc_sources}

For an exact example, let hidden borrower type $U$ be Bernoulli with probability
$1/2$, let observed treatment $A$ agree with $U$ with probability $3/4$, and let
the outcome be $Y=2A+4U$. Then
\[
 E[Y\mid A=1]=5,\qquad E[Y\mid A=0]=1,
 \qquad E[Y\mid\mathrm{do}(A=1)]-E[Y\mid\mathrm{do}(A=0)]=2.
\]
The association is four, while the intervention effect is two because changing
$A$ leaves the population law of $U$ unchanged. At a stipulated value of
USD 100 per outcome unit and action cost USD 250, the association-based decision
reports net USD 150 but the causal calculation gives USD $-50$. The numerical
example is an authored finite-law calculation; it is neither a package
benchmark nor estimated bank economics.
\citep{exec23_pymc_reference}

Even a perfect posterior computation cannot resolve every causal ambiguity.
The structural models $A=U,Y=A$ and $A=U,Y=U$, with a fair binary $U$, induce
the same observed law for $(A,Y)$ but treatment effects one and zero.
Under equal observational likelihoods, posterior odds between these models
remain their prior odds. A narrower posterior obtained by choosing one model
is conditional on that choice, not evidence that observations distinguished
the models. A useful separate computational oracle is the conjugate update
from a uniform beta prior after six successes and four failures: the posterior
is beta$(7,5)$ with mean $7/12$ and variance $35/1872$. Agreement with this
oracle checks a declared probability calculation; it does not validate a
causal interpretation of those successes.
\citep{exec23_pymc_reference}

The transformation API also allows interventions that reference the original
variable. The inspected code warns and gives the replacement a new name with
a \texttt{do\_} prefix. Such an expression needs an explicit scientific
interpretation, for example a declared shift policy; it must not be silently
described as a fixed assignment. Optional pruning removes variables detached
from observed variables, and its helper explicitly refuses models containing
Potentials because their prior-versus-likelihood role is ambiguous. Therefore
pruning cannot substitute for reviewing the target and retained information.
\citep{exec23_pymc_sources}

\subsection{Stored draws, regenerated variables and changed expressions}
\label{exec23:pymc-prediction}

The selected forward-sampling implementation separates output selection from
regeneration. \texttt{var\_names} controls what is returned; it does not request
a fresh draw. \texttt{sample\_vars} identifies variables to regenerate, while
\texttt{freeze\_vars} preserves trace values. Random variables present in the
trace and not explicitly selected for regeneration are automatically frozen.
Missing random variables, including ordinary observed outcomes absent from
the trace, are regenerated. The code detects some changed-ancestor cases and
issues an implicit-freeze warning, but warning availability does not establish
that the chosen predictive experiment is appropriate.
\citep{exec23_pymc_sources}

Suppose a stored draw contains $X=1,Y=1$ from a model with
$Y\mid X\sim N(X,1)$. After changing $X$ to three, retaining the stored $Y$
returns one, whereas regenerating $Y$ from the changed model has conditional
mean three. Both operations can be intentional, but they answer different
questions. A policy simulation must identify which parameter draws remain
conditioned on the original evidence and which post-intervention variables
are generated anew. Otherwise an apparently successful transformation can
leave precisely the outcome of interest tied to the old experiment.
\citep{exec23_pymc_sources,exec23_pymc_reference}

Changed deterministic expressions require their own attention. The inspected
documentation explicitly describes replacing a deterministic expression without
changing any random variable or data value: the predictive machinery can then
reuse its stored value unless regeneration is requested. For two illustrative
trace points $X=1,2$, replacing $X^2$ by $X^3$ changes the desired values from
$(1,4)$ to $(1,8)$ and their mean from $5/2$ to $9/2$. A model digest must
therefore identify expressions, not just variable names and current data.
Name-based trace matching with compatible shapes and coordinates is a software
matching rule, not proof of causal or semantic equivalence between models.
\citep{exec23_pymc_sources,exec23_pymc_reference}

Predictive handling of Potentials adds another boundary. The inspected prior
and posterior predictive routines warn that their direct effect on other
parameters is ignored during forward sampling. A fair binary generator with
a hard weight excluding zero has weighted probability one of a positive value,
whereas the unconstrained forward generator still has probability $1/2$.
This witness explains why a density factor cannot automatically be treated
as a generative rule. It does not imply that all effects of a Potential vanish:
posterior ancestors supplied by the trace may already reflect that factor.
The worker must identify the factor's role, whether it enters the inferential
posterior, and how the requested predictive law is constructed.
\citep{exec23_pymc_sources,exec23_pymc_reference}

The selected upstream tests include intervention, self-reference, dimensional
and predictive cases, as well as a constant-intervention predictive case marked
as an expected failure. These are useful conformance leads, not executed results
in this project. The proposed fixture set must exercise both shared and constant
interventions and inspect the actual result and warning, without converting
an upstream test marker into a claim that failure was reproduced here.
\citep{exec23_pymc_sources}

\subsection{Computation and diagnostics must remain visible to the caller}
\label{exec23:pymc-decision}

Sampler identity cannot safely be left implicit. At the selected revision,
the default dispatch can choose nutpie when the installed environment, model,
compilation mode and arguments permit it, otherwise the PyMC sampler. The
external paths impose restrictions on variable types and supplied steps,
and reject custom trace backends and per-draw callbacks. Consequently,
partial-trace persistence described for a PyMC trace backend cannot be
advertised for every sampler choice. Persisting draws also does not by itself
establish resumable chain state, independent restarts or distributed recovery.
\citep{exec23_pymc_sources}

The inspected convergence routine reports missing posterior data, too few
draws, a single chain and fewer than four chains separately. It checks
reported divergences and tree-depth statistics; missing statistics return no
corresponding warning. Its inspected thresholds include an $\widehat R$
warning above 1.01 and an effective-sample-size warning below 100 per chain.
These are the package's conditional diagnostics, not a theorem of convergence
or a guarantee that an intervention model is correct. NUTS source documentation
also describes energy, trajectory-depth and timing statistics and distinguishes
CPU from wall time. The product must preserve missing, unavailable, adverse
and satisfactory diagnostics as different states, rather than flatten them
into a single successful API status.
\citep{exec23_pymc_sources}

The current disposition is a source-qualified reference with a conditional
native Python worker route. A blanket port to other languages would reproduce
graph rewriting and predictive rules at considerable maintenance cost without
a demonstrated advantage. A concrete prototype should instead pin the release
or commit, Python, PyTensor, compiler, sampler and dependencies; register exact
conjugate and finite latent-model oracles; exercise intervention/freeze/sample
semantics, Potentials, dimensions and warnings; and measure compilation,
effective samples, memory, predictive cost and recovery together. Scientific
owners must justify model and intervention meaning, while maintainers own
conversion conformance and operators own budgets and persistence. Library,
CLI and MCP should expose the same target and refusal conditions. LLM-generated
models can be proposed for review, but cannot certify their own causal graph,
identification assumptions or scientific validity. Native qualification,
general causal identification and bank adoption remain open.
\citep{exec23_pymc_sources,exec23_pymc_reference}
